
\section{R1CS with scheduled witnesses. \tbc}
Two-round R1CS, to support lookups into tables, which are baked into the R1CS matrices.
The witness $\vec w$ consists of 
\begin{enumerate}
    \item 
    Round-1 witness $\vec w_1$: 
    These witnesses are independent of the verifier challenge
    Includes public inputs, constants, and lookup table multiplicities.
    \item 
    Round-2 witness $\vec w_2$: 
    These witnesses depend on the verifier challenges.
    These are the verifier challenges themselves, the logUp inverses and their sums.
\end{enumerate}
\ulrich{Let us elaborate on the usefulness of 2-stage R1CS beyond lookups.}
\ulrich{%
Need to define $\mathscr R_{main}$ and $\mathscr R_{aux}$.
Target $\vec w_1\in \mathscr R_{main} \cap \mathscr R_{aux}$
}
% The verifier challenges for logUp are dedicated inputs of the circuit, and the matrices $A$, $B$, $C$ cover the logUp logic. 
The complete witness for the constraint system is of the form 
\begin{equation*}
%\label{e:witness:split}
\vec w = \vec w_1\| \vec w_2,
\end{equation*}
with $\vec w_1\in \mathbb F_q^{N_1}$ and $\vec w_2\in \mathbb F_q^{N_2}$, 
\ulrich{Concatenation is never used anywhere.}
\ulrich{%
It seems more efficient to treat logup with several basefield challenges.
This is for sure the case for the witness builder.
Memory constraints in the zerocheck (due to the higher-dimensional hypercube) can be overcome by univariate skip.
}
subject to $M\geq 1$ constraints
\begin{equation*}
% \label{e:R1CS}
(A\cdot \vec w) \odot (B\cdot \vec w) = C\cdot \vec w,
\end{equation*}
with $A, B, C\in \mathbb F^{K\times N}$, and where $\odot$ is the entrywise (Hadamard) product.
As for the witness, we split the matrices $M= A, B, C$ into 
\begin{equation*}
M = 
\begin{pmatrix}
M_1 | M_2
\end{pmatrix}, 
\end{equation*}
with $M_1\in \mathbb F_q^{K \times N_1}$ and $M_2\in \mathbb F_q^{K\times N_2}$ corresponding to the $\vec w_1$- and $\vec w_2$-parts. 
% \ulrich{Would $M_2$ over $F$ make sense?}
Typically the system is partially decoupled, with each $A_2, B_2, C_2$ starting with a large block of zero rows, but we currently do not leverage simpler representations than the ones above. 
% \ulrich{Do we?}

% \[
% \left[
% \left(
% \begin{array}{c|c}
% A_{1,1}  & 0
% \\\hline
% A_{2,1} & A_{2,2} 
% \end{array}
% \right) 
% \cdot
% \begin{pmatrix}
% \vec w_1 \\ \vec w_2
% \end{pmatrix}
% \right]
% \odot
% \left[
% \left(
% \begin{array}{c|c}
% B_{1,1}  & 0
% \\\hline
% B_{2,1} & B_{2,2} 
% \end{array}
% \right) 
% \cdot
% \begin{pmatrix}
% \vec w_1 \\ \vec w_2
% \end{pmatrix}
% \right]
% =
% \left(
% \begin{array}{c|c}
% C_{1,1}  & 0
% \\\hline
% C_{2,1} & C_{2,2} 
% \end{array}
% \right) 
% \cdot
% \begin{pmatrix}
% \vec w_1 \\ \vec w_2
% \end{pmatrix}
% \]
% \[
% M = 
% \left(
% \begin{array}{c|c}
% M_{1,1}  & 0
% \\\hline
% M_{2,1} & M_{2,2} 
% \end{array}
% \right) 
% \]
% with $M_{2,1}$ and $M_{2,2}$ very sparse (most of the rows have at most 2 entries).

We need:
\begin{itemize}
    \item 
    Definition of the first stage R1CS with auxiliary lookup relation $\mathscr R_\text{R1CS} \cap \mathscr R_\text{aux}$
    \item
    Definition of randomized R1CS (with two stages), and its soundness error
\end{itemize}

\subsubsection{Split-witness builder.}
There is a split-witness builder for stepwise computation of the complete witness.